\pdfminorversion=4
\documentclass[aspectratio=169]{beamer}

\mode<presentation>
{
  \usetheme{default}
  \usecolortheme{default}
  \usefonttheme{default}
  \setbeamertemplate{navigation symbols}{}
  \setbeamertemplate{caption}[numbered]
  \setbeamertemplate{footline}[frame number]  % or "page number"
  \setbeamercolor{frametitle}{fg=white}
  \setbeamercolor{footline}{fg=black}
} 

\usepackage[english]{babel}
\usepackage[utf8x]{inputenc}
\usepackage{tikz}
\usepackage{courier}
\usepackage{array}
\usepackage{bold-extra}
\usepackage{minted}
\usepackage[thicklines]{cancel}
\usepackage{fancyvrb}
\usepackage{ulem}

\xdefinecolor{dianablue}{rgb}{0.18,0.24,0.31}
\xdefinecolor{darkblue}{rgb}{0.1,0.1,0.7}
\xdefinecolor{darkgreen}{rgb}{0,0.5,0}
\xdefinecolor{darkgrey}{rgb}{0.35,0.35,0.35}
\xdefinecolor{darkorange}{rgb}{0.8,0.5,0}
\xdefinecolor{darkred}{rgb}{0.7,0,0}
\definecolor{darkgreen}{rgb}{0,0.6,0}
\definecolor{mauve}{rgb}{0.58,0,0.82}

\title[2021-02-03-pyhep-numba-tutorial]{{\vspace{0.4 cm}\\ \LARGE PyHEP Numba Tutorial} \vspace{0.2 cm}\\ or \vspace{0.2 cm}\\ Getting Numba to Do What You Want}
\author{Jim Pivarski}
\institute{Princeton University -- IRIS-HEP}
\date{February 3, 2021}

\usetikzlibrary{shapes.callouts}

\begin{document}

\logo{\pgfputat{\pgfxy(0.11, 7.4)}{\pgfbox[right,base]{\tikz{\filldraw[fill=dianablue, draw=none] (0 cm, 0 cm) rectangle (50 cm, 1 cm);}\mbox{\hspace{-8 cm}\includegraphics[height=1 cm]{princeton-logo-long.png}\hspace{0.1 cm}\raisebox{0.1 cm}{\includegraphics[height=0.8 cm]{iris-hep-logo-long.png}}\hspace{0.1 cm}}}}}

\begin{frame}
  \titlepage
\end{frame}

\logo{\pgfputat{\pgfxy(0.11, 7.4)}{\pgfbox[right,base]{\tikz{\filldraw[fill=dianablue, draw=none] (0 cm, 0 cm) rectangle (50 cm, 1 cm);}\mbox{\hspace{-8 cm}\includegraphics[height=1 cm]{princeton-logo.png}\hspace{0.1 cm}\raisebox{0.1 cm}{\includegraphics[height=0.8 cm]{iris-hep-logo.png}}\hspace{0.1 cm}}}}}

% Uncomment these lines for an automatically generated outline.
%\begin{frame}{Outline}
%  \tableofcontents
%\end{frame}

% START START START START START START START START START START START START START

\begin{frame}{To optimize or not to optimize?}
\Large
\vspace{0.4 cm}
\textcolor{darkblue}{Scientific computing is of two minds about accelerating code:}

\vspace{0.2 cm}
\begin{center}
\begin{minipage}{0.75\linewidth}
\begin{block}{}
``Tweaking performance should take a back seat to the science.''
\end{block}

\begin{block}{}
``If it doesn't run within time/memory constraints, we can't perform this measurement.''
\end{block}
\end{minipage}
\end{center}

\vspace{0.5 cm}
\uncover<2->{Both are true, so it's a balance. The objective: \textcolor{darkorange}{\bf ``time to insight.''}}
\end{frame}

\begin{frame}{Compilation}
\Large
\vspace{0.5 cm}
Roughly speaking, compilation makes code faster.

\vspace{0.5 cm}
\uncover<2->{\textcolor{darkblue}{More precisely, removing hurdles takes the slowness out of code.}}
\large

\vspace{0.5 cm}
\begin{uncoverenv}<2->
\hfill\begin{minipage}{0.97\linewidth}
\begin{itemize}
\item Avoid type-checking every value in a loop over identically typed data.
\item Avoid following pointers from objects to objects.
\item Avoid nonsequential memory access.
\item Avoid translating virtual machine \\ instructions into real instructions.
\item Avoid garbage collector searches \\ for many small objects.
\end{itemize}
\end{minipage}

\vspace{-2.5 cm}
\hfill \includegraphics[width=5.5 cm]{hurdle9.jpg}
\end{uncoverenv}
\end{frame}

\begin{frame}[fragile]{Compilation is not a monolithic thing}
\Large

\vspace{0.5 cm}
Though usually thought of as single step applying all optimizations,

\vspace{0.2 cm}
\hspace{0.6 cm} \begin{minipage}{0.9\linewidth}
\begin{minted}{bash}
c++ my_source_code.cpp -o now_its_faster
\end{minted}
\end{minipage}

\vspace{0.25 cm}
compilation can be piecemeal and applied at any time.

\vspace{0.5 cm}

\mbox{ } \hfill \includegraphics[height=5 cm]{lisp-faster-than-c.png} \hfill \mbox{ }
\end{frame}

\begin{frame}{Compilers are everywhere}
\large
\vspace{0.35 cm}

\begin{itemize}\setlength{\itemsep}{0.27 cm}
\item Python is ``compiled'' from source code to numeric bytecodes that run in a virtual machine. (At least it isn't scanning the source code over and over!)
\item \mintinline{python}{re.compile("...")} $\longrightarrow$ regular expression to finite state machine.
\item \mintinline{python}{struct.Struct("...")} $\longrightarrow$ same thing for parsing bytes.
\item \mintinline{python}{numexpr.evaluate("...")} $\longrightarrow$ NumExpr's virtual machine for evaluating math formulas is faster than NumPy.
\item \mintinline{python}{ROOT.gInterpreter.Declare("...")} $\longrightarrow$ ROOT's Cling compiles strings of C++ code to machine-native bytecode instructions.
\item \mintinline{python}{rdf.Define("...")} $\longrightarrow$ ROOT's RDataFrame also uses Cling.
\item Numba and PyPy compile Python to machine-native bytecode instructions.
\item Julia is a language designed around this \only<1>{just-in-time (JIT) compilation.}\only<2>{\textcolor{darkorange}{\bf just-in-time (JIT) compilation.}}
\end{itemize}
\end{frame}

\begin{frame}{Just-in-time (JIT) compilers for Python}
\large

\begin{columns}
\column{0.5\linewidth}
\begin{center}
\includegraphics[width=0.75\linewidth]{numba-logo.png}
\end{center}

Runs within normal CPython (can be used with other Python libraries).

\vspace{0.2 cm}

Only speeds up functions labeled with the \mintinline{python}{@numba.jit} decorator.

\vspace{0.2 cm}

\only<1>{Only a subset of Python \mbox{features\hspace{0.75 cm}} and functions can be JIT-compiled.}\only<2>{\textcolor{darkorange}{\bf Only a subset of Python features and functions can be JIT-compiled.}}

\vspace{0.2 cm}

Gains are often \textcolor{darkblue}{factors of 100's}.

\column{0.5\linewidth}
\begin{center}
\includegraphics[width=0.75\linewidth]{pypy-logo.jpg}
\end{center}

Replaces the CPython process; not all libraries/versions of Python work.

\vspace{0.2 cm}

Speeds up the whole program. \mbox{\hspace{4 cm}}

\vspace{0.2 cm}

All Python features are JIT-compiled. \mbox{\hspace{4 cm}}

\vspace{0.2 cm}

Gains are often \textcolor{darkblue}{factors of several}.

\end{columns}
\end{frame}




\end{document}
